\chapter{Flows} % chapter* je necislovana kapitola

In this chapter, we aim to introduce all concepts and current knowledge about flows and related graph theory terms. Mainly, we explain \emph{Chebyshev nowhere-zero flows}, which
have arisen in the recent preprint coauthored by the author \cite{preprint}. Moreover, we present a \emph{flow-pair} -- a sufficient condition for a graph to have
some nowhere-zero flow and some Chebyshev flow, also stated in that paper. We also establish a notion of \emph{rich flows} not allowing same flow values on darts around a vertex.

\section{Group flows}

Intuitively, flows represent the dynamics of a fluid in some network. Consider that $x$ gallons of water from a vertex $u$ flow to a vertex $v$. When we ask, how many gallons of water from a vertex $v$ flow to a vertex $u$, we should say $-x$. Hence, when working with flows, we naturally need inverses. Moreover, if the system is stable, the amount of water that flows in should be the same as the amount of water that drains out. This can be reformulated as ``(possibly negative) amounts of water flowing into a vertex sum up to zero''. This motivates the need for a neutral element. So, we end up with a group. Therefore, the original definition of nowhere-zero flows uses a group-theory setting with additive notation.

\begin{definition}
Let $G=(V, E)$ be a graph, $\mathcal A$ be an abelian group, and $F\subseteq\mathcal A$ be the set of permitted flow values closed under inverse. An \emph{$(\mathcal A, F)$-flow} on $G$ is an assignment $\varphi\colon D(G)\to F$ of flow values to the darts satisfying
$\varphi(\vv{uv})=-\varphi(\vv{vu})$ for any dart $\vv{uv}\in D(V)$ and
\begin{equation}
\sum_{d\in D(\{v\})} \mkern-10mu\varphi\left(d\right) = 0 \label{eq:conservation}
\end{equation}
for each vertex $v$ of $G$. The equality \eqref{eq:conservation} is often referred to as the \emph{flow conservation constraint}. If $F=\mathcal A$, we use an abbreviation \emph{$\mathcal A$-flow}, and if $F=\mathcal A\setminus{0}$, we use the term \emph{nowhere-zero $\mathcal A$-flow} (or abbreviated \emph{NZ $\mathcal A$-flow}).
\end{definition}

Observe that the nowhere-zero condition is crucial in avoiding trivial flows, since otherwise, we are able to set all flow values to zero, and it is a valid flow. The flow conservation constraint \eqref{eq:conservation} is generalisable to any subset of vertices. Hence, having a graph $G$ with a bridge $e$, the interpretation of the flow conservation constraint for one of the components of $G - e$ results in the fact, that $G$ admits no NZ flow. Therefore, nowhere-zero flows are researched only on bridgeless graphs. However, by an interesting argument, the number of NZ $\mathcal A$-flows of some graph depends only on the order of the group.

\begin{theorem} \emph{\cite[p. 82]{tutte}}
Let $\mathcal A_1, \mathcal A_2$ be two finite abelian groups of the same order and $G$ be a graph. Then $G$ admits an NZ $\mathcal A_1$-flow if and only if it admits an NZ $\mathcal A_2$-flow.
\end{theorem}

\section{Integral flows}

Results in the previous section motivated, that the specific group is not relevant, relevant is only its order. We can use this to get out of the group setting and replace its elements by integers.

\begin{definition}
Let $k\geq 2$ be a positive integer. A \emph{$k$-flow} is a $(\mathbb Z, \{0,\pm 1, \pm 2, \dots, \pm (k-1)\})$-flow. Moreover, a \emph{nowhere-zero $k$-flow} (or \emph{NZ $k$-flow}) is a $(\mathbb Z, \{\pm 1, \pm 2, \dots, \pm (k-1)\})$-flow.
\end{definition}

Since integers are well-ordered, we also establish a notation $\varphi(e)$ for $e=uv\in E$, where $\varphi(e) = \max\{\varphi(\vv{uv}), \varphi(\vv{vu})\}$. Hence, we can also take a \emph{positive} flow -- for each edge, we take a dart with the positive flow value. Now, to see that this newly defined notion truly replaces the problem in the group theory setting, we need the following theorem.

\begin{theorem} \emph{\cite[p. 207]{jaeger1979flows}}
For each $\mathbb Z_k$-flow $\varphi$ there exists a $k$-flow $\varphi'$ satisfying $\varphi'\left(d\right)\equiv\varphi\left(d\right)\pmod{k}$. \label{th:finite_group_flows}
\end{theorem}

\begin{remark}
Since the preceding theorem preserves modulo, it also preserves zeroes, and hence the same holds for NZ flows.
\end{remark}

It is not hard to see that a graph with an NZ $k$-flow also admits an NZ $K$-flow for any $K>k$. Hence, it makes sense to seek the smallest $k$ such that a graph admits an NZ $k$-flow.

\begin{definition}
Let $G$ be a bridgeless graph. A \emph{flow number of $G$} is
\begin{equation*}
\Phi(G) := \min\{k\mid\exists\text{NZ } k\text{-flow on }G\}.\label{eq:flow_number}
\end{equation*}
\end{definition}

However, it turns out that this rate of complexity achieves only a few values.

\begin{theorem}[$6$-flow theorem] \emph{\cite[p. 133]{seymour}}
There exists an NZ $6$-flow on any bridgeless graph.\label{th:6_flow}
\end{theorem}

Seymour constructed this NZ $6$-flow from two simpler flows from the fact, that if $a\in\{0,\pm 1\}$, $b\in\{0,\pm 1,\pm 2\}$ and not both $a, b$ are zero, then $3a+b\in\{\pm 1, \pm 2, \pm 3, \pm 4, \pm 5\}$.

\begin{lemma} \emph{\cite[p. 132]{seymour}}
For any bridgeless graph, there exist a $(\mathbb Z\times\mathbb Z, (\{0,\pm 1\}\times\{0, \pm 1, \pm 2\})\setminus\{(0, 0)\})$-flow. \label{lem:2_flow_3_flow_seymour}
\end{lemma}

On the other hand, there is no known graph with $\Phi(G)=6$, all known graphs have $\Phi(G)\leq 5$. Therefore, the following conjecture about the NZ $5$-flow by Tutte still holds unresolved.

\begin{conjecture}[$5$-flow conjecture] \emph{\cite[p. 83]{tutte}}
There exists an NZ $5$-flow on any bridgeless graph.\label{conj:5_flow}
\end{conjecture}

As for many problems in graph theory, it is sufficient to explore the problem only on a smaller subset of graphs. There exists a folklore reduction of nowhere-zero flows to cubic graphs.
\begin{proposition}
There exists an NZ $k$-flow on any bridgeless graph if and only if there exists an NZ $k$-flow on any cubic bridgeless graph.\label{prop:reduction}
\end{proposition}

So if there is a counterexample to Tutte's $5$-flow conjecture, there is also a cubic counterexample. Moreover, in terms of finding such counterexample, we can also avoid colourable ones and limit research to snarks, as a corollary of the next proposition.

\begin{proposition} \emph{\cite[pp. 160, 161]{diestel}}
A cubic graph has an NZ $4$-flow if and only if it is $3$-edge-colourable.
\end{proposition}

There are more constraints regarding the smallest counterexample of Conjecture \ref{conj:5_flow}. One of the first results concerns snarks of small oddness.

\begin{theorem} \emph{\cite[p. 89]{nearly_f4}}
Let $G$ be a snark of oddness $2$. Then, $\Phi(G)\leq 5$. \label{th:oddness_1D}
\end{theorem}

As a consequence, since each weak snark has oddness $2$ (which is considered to be folklore), this also covers all weak snarks. Hence, any counterexample to $5$-flow conjecture must be a strong snark of oddness at least $4$. Moreover, we know some properties of the potential smallest counterexample to the $5$-flow conjecture.

\begin{theorem} \emph{\cite[p. 142]{min_ctrex_cyc_edge_con}}
If the $5$-flow conjecture is false, the smallest counterexample is a cyclically $6$-edge-connected snark.
\end{theorem}

\begin{theorem} \emph{\cite[p. 364]{oddness_4}}
If the $5$-flow conjecture is false, the smallest counterexample has oddness at least $6$.
\end{theorem}

\begin{theorem} \emph{\cite[p. 15]{min_ctrex_girth}}
If the $5$-flow conjecture is false, the smallest counterexample has girth at least $12$.
\end{theorem}

\section{Chebyshev flows}

There were some attempts to generalise nowhere-zero flows to real numbers and more dimensions. By generalising integer nowhere-zero flows to real numbers, researchers found out that restricting only zero flow value is not sufficient, since any flow consisting of non-zero flow values can be made ``smaller'' and then, the notion of flow number disappears. Therefore, we need to require an interval of forbidden flow values, and hence, one-dimensional real nowhere-zero flows are allowed to use flow values from $[1-r,r-1]\setminus(-1, 1)$ \cite{goddyn_fractional_nzf}. The generalisation of real flows to more dimensions \cite{euclidean_flows} picks only vectors with the Euclidean norm from the mentioned interval $[1,r-1]$. It works, but it did not have such nice properties, as we would expect. However, changing a norm turned out to be a good idea \cite{preprint}, so we will consider only Chebyshev multidimensional nowhere-zero flows.

\begin{definition}
A \emph{$d$-dimensional open and closed balls of radii $r$ with respect to Chebyshev norm} are
\begin{align*}
B_d^\infty(r) &= \left\{x\in\mathbb R^d~\middle|~\|x\|_\infty < r\right\}\text{, and}\\
\overline{B_d^\infty}(r) &= \left\{x\in\mathbb R^d~\middle|~\|x\|_\infty \leq r\right\},
\end{align*}
respectively.
\end{definition}

\begin{definition}
Let $d$ be a positive integer and $r\geq 2$ be a real constant. A \emph{$d$-dimensional Chebyshev nowhere-zero $r$-flow} (or \emph{$d$D ChNZ $r$-flow}) is a $\left(\mathbb R^d, \overline{B_d^\infty}(r-1)\setminus B_d^\infty(1)\right)$-flow.
\end{definition}

Now, for $e=uv$, the notation $\varphi(e)$ denotes the lexicografically larger one from $\varphi\left(\vv{uv}\right), \varphi\left(\vv{vu}\right)$. The flow number can be defined in a similar way to integral nowhere-zero flows. The only problem is that if we take a set $\{r\mid\exists d\text{D ChNZ }r\text{-flow on }G\}$, it may not have a minimum, so a formally correct definition would use infimum instead of minimum. However, the infimum of the flow number candidates is always feasible \cite[p. 101]{euclidean_flows}, so we can define the flow number straightforwardly with the minimum.

\begin{definition}
Let $G$ be a bridgeless graph and $d$ be a positive integer. A \emph{$d$-dimensional Chebyshev flow number} of $G$ is
\begin{equation*}
\Phi_d^\infty(G) := \min\{r\mid\exists d\text{D ChNZ }r\text{-flow on }G\}.
\end{equation*}
\end{definition}

One special case worth mentioning is a \emph{circular NZ flow} – a $1$D ChNZ flow. In one dimension, it does not matter, which norm is used to constrain permitted flow values, so we can omit it from the naming. Similarly, $\Phi_1^\infty$ can be replaced by $\Phi_1$. As we show in the following proposition, there is a relation between circular flow number and integral flow number of a graph – namely $\Phi(G) = \lceil \Phi_1(G)\rceil$.

It has been proved \cite[p. 4]{preprint} that the Chebyshev flow number is always rational. Moreover, the authors found a normal form of Chebyshev flows, as stated in the following proposition.% Since we need the explicit construction in the latter proofs, we will give a sketch of this proof.

\begin{proposition} \emph{\cite[p. 6]{preprint}}
Let $G$ be a graph with a ChNZ $p/q$-flow $\varphi$ for some positive integers $p, q$. Then, there exists a $d$D ChNZ $p/q$-flow such that every coordinate of every flow value is one of the two nearest integral multiples of $1/q$ to the corresponding coordinate in $\varphi$. \label{prop:chebyshev_normal_form}
\end{proposition}

%\begin{proof}[Sketch of the proof]
% If there is a dart with flow value, whose $i^{\text{th}}$ coordinate is not an integral multiple of $1/q$, we can find an oriented cycle $\vv{\mathcal{C}}$ containing that dart with a property that no dart of this cycle has flow value with $i^{\text{th}}$, which is an integral multiple of $1/q$. Otherwise, in some vertex, all but one dart have $i^{\text{th}}$ coordinate of flow value, which is a multiple of $1/q$, and hence the sum of values on all darts (which is zero by the flow conservation constraint) is not a multiple of $1/q$, a contradiction. Now consider for each $d\in\vv{\mathcal{C}}$ the maximal possible adjustment of $\varphi\left(d\right)_i$ such that its absolute value remains in the range $\left[1,p/q-1\right]$. Formally,
% \begin{equation*}
%  \delta\left(d\right) = \begin{cases}
%   p/q-1-\varphi\left(d\right)_i & \text{if } \varphi\left(d\right)_i > 0,\
%   -1-\varphi\left(d\right)i & \text{if } \varphi\left(d\right)i < 0.
%  \end{cases}
% \end{equation*}
% Clearly $\delta\left(d\right) > 0$ for any dart on cycle $\vv{\mathcal C}$, since both $p/q-1, -1$ are integral multiples of $1/q$. Thus, by sending $\delta{\text{min}}=\min{d\in\vv{\mathcal{C}}} \delta\left(d\right)$ along $i^{\text{th}}$ coordinate, the number of conflicting coordinates decreases. By repeating this process, we get to a desired normal form.
%\end{proof}

Now we recall Seymour's flows from Lemma \ref{lem:2_flow_3_flow_seymour} – a $2$-flow and a $3$-flow with a property that no edge has both flow values equal to zero. By simply using the value of the $2$-flow as the first coordinate and the value of the $3$-flow as the second coordinate, we get a Chebyshev nowhere-zero flow. Moreover, this flow gives us a nice bound on the Chebyshev flow number.

\begin{proposition} \emph{\cite[p. 8]{preprint}}
Each bridgeless graph admits a 2D ChNZ $3$-flow.\label{prop:chebyshev_upper_seymour}
\end{proposition}

Again, no graph with $\Phi_2^\infty(G)=3$ is known. Furthermore, there is no known graph with $\Phi_2^\infty(G)>5/2$. Therefore, the next conjecture is interesting to examine.

\begin{conjecture} \emph{\cite[p. 22]{preprint}}
Each bridgeless graph admits a 2D ChNZ $5/2$-flow.\label{conj:chebyshev_upper_seymour}
\end{conjecture}

Similarly to integral flows, the problem can be reduced to cubic graphs and then to snarks, since colourable cubic graphs have a small flow number.

\begin{theorem} \emph{\cite[p. 14]{preprint}}
A cubic graph has a 2D ChNZ $2$-flow if and only if it is $3$-edge-colourable.\label{th:2_chnzf_iff_3_col}
\end{theorem}

\section{Flow-pairs}

As an NZ $6$-flow and a 2D ChNZ $3$-flow can be constructed from two simpler flows, something similar can be made for smaller flow numbers with two other flows \cite[p. 19]{preprint}.

\begin{definition}
Let $p\leq q$ be positive integers. A \emph{$p/q$-flow-pair} of a graph $G$ is a $(\mathbb Z\times\mathbb Z, F)$-flow with
\begin{equation*}
F = \{(0, f)\mid q\leq|f|\leq p+q\}\cup\{(1,f)\mid |f|\leq p+q\}.
\end{equation*}
\end{definition}

\begin{remark}
The structure in Lemma \ref{lem:2_flow_3_flow_seymour} is a $1/1$-flow-pair.
\end{remark}

One possible ambiguity of this definition is that since $p/q$ is written as a fraction, the $2/6$-flow-pair can be mixed up with the $1/3$-flow-pair. However, their existence is equivalent (by a similar argument to Proposition \ref{prop:chebyshev_normal_form}), so we can WLOG assume that $p,q$ are relatively prime.

\begin{lemma} \emph{\cite[p. 19]{preprint}}
A bridgeless graph $G$ with a $dp/dq$-flow-pair also has a $p/q$-flow-pair and vice versa.
\end{lemma}

We mentioned that the $p/q$-flow-pair is a sufficient condition for the existence of some nowhere-zero flows. Those flows are somehow related to the value $p/q$, and that is the reason to write it as a fraction in the definition of the flow-pair.

\begin{proposition} \emph{\cite[pp. 19, 20]{preprint}}
Consider a bridgeless graph $G$ with a $p/q$-flow-pair $\varphi=(\varphi_1, \varphi_2)$. Then, there are a 2D ChNZ $(2+p/q)$-flow $\varphi'=(\varphi_1, \varphi_2/q)$ and a circular NZ $(4+2p/q)$-flow $\varphi''=(2+p/q)\varphi_1+\varphi_2/q$ on $G$.\label{prop:nzf_from_flow_pair}
\end{proposition}

Note that if there is a graph, which is a counterexample to Tutte's 5-flow conjecture, it should not have the $1/2$-flow-pair. Hence, it would be interesting to classify graphs with $1/2$-flow-pair. By a computer, it was checked that this flow-pair exists for snarks up to $34$ vertices, so it looks like the $1/2$-flow-pair exists always.

\begin{conjecture} \emph{\cite[p. 22]{preprint}}
Each bridgeless graph has a $1/2$-flow-pair.
\end{conjecture}

In general, edges with a non-zero value of the first coordinate form a set of cycles in $G$. In cubic graphs, these cycles should be clearly vertex-disjoint. Moreover, if considering other than the trivial $1/1$-flow-pair, these cycles together cover all vertices.

\begin{lemma} \emph{\cite[p. 21]{preprint}}
Let $p<q$ be positive integers. Consider a cubic graph $G$ with a $p/q$-flow-pair $\varphi=(\varphi_1, \varphi_2)$. Then the set of edges $e$ with a property $\varphi_1(e)\neq 0$ is a $2$-factor of $G$.\label{lem:flow_pair_2_factor}
\end{lemma}

\section{Rich flows}

Recently, in the research of flows, there was an idea to forbid two darts surrounding a vertex to have flow values with the same absolute value \cite{rich_flows, rich_flows_not_cubic}. The motivation is that if we have two darts from the same vertex with $\varphi\left(d_1\right)+\varphi\left(d_2\right)=0$ (a so-called \emph{confluent pair}), we can split the considered vertex of the degree $d$ to vertices of degrees $2$ and $d-2$ without influencing flow properties of a graph. On the other hand, a so-called \emph{contrafluent pair}, where $\varphi\left(d_1\right)=\varphi\left(d_2\right)$, requires that the set of flow values contains some element and its double. However, we cannot find such a pair in a $p/q$-flow-pair with $p<q$.

\begin{definition}
An NZ $k$-flow $\varphi$ is called \emph{rich} if for each pair of adjacent darts $d_1$, $d_2$ it holds that $\left|\varphi\left(d_1\right)\right|\neq\left|\varphi\left(d_2\right)\right|$.
\end{definition}

We already know that a graph admits a nowhere-zero flow if and only if it is bridgeless. In the theory of rich flows, there is one more obstruction – a $2$-edge-cut with the common vertex. It can be shown that these two conditions are also sufficient for a graph to have the rich nowhere-zero flow. We call such graphs \emph{rich flow admissible}.

\begin{definition}
Let $G$ be a rich flow admissible graph. A \emph{rich flow number of $G$} is
\begin{equation*}
R(G) := \min\{k\mid\exists\text{rich NZ } k\text{-flow on }G\}.
\end{equation*}
\end{definition}

Unlike group, integral, circular, and multidimensional flows, rich flows and bounding the rich flow number cannot be reduced to the class of cubic graphs. The simple reason is that we clearly need at least $\Delta(G)$ different flow values and therefore $R(G)\geq\Delta(G)+1$. On the other hand, there is a linear upper bound depending only on $\Delta(G)$ \cite{rich_flows_not_cubic}.

However, if we limit our view only to the cubic graphs, we get similar results as for integral nowhere-zero flows. Again, it suffices to examine the snarks.
\begin{proposition} \emph{\cite[p. 31]{rich_flows}}
A cubic graph has a rich NZ $5$-flow if and only if it is $3$-edge-colourable.
\end{proposition}

And from the computational results, the gap between the colourable and uncolourable cubic graphs is again equal to $1$. However, with the proven results, we are still far from the desired value $6$.
\begin{conjecture} \emph{\cite[p. 33]{rich_flows}}
There exists a rich NZ $6$-flow on any bridgeless cubic graph.
\end{conjecture}

\begin{proposition} \emph{\cite[p. 33]{rich_flows}}
There exists a rich NZ $11$-flow on any bridgeless cubic graph.
\end{proposition}

For snarks that are ``nearly colourable'', we have a result similar to Corollary \ref{th:oddness_1D} –- we can reduce the size of the gap from $5$ to $3$.
\begin{proposition} \emph{\cite[p. 40]{rich_flows}}
Let $G$ be a weak snark. Then, $R(G)\leq 9$. \label{prop:weak_rich_flow}
\end{proposition}
